% =====================================================================
% A2 — Section 1: Introduction.  DRAFT 1, 6 Oct 2026.
% Problem -> what the paper claims -> why this claim -> what we did ->
% what we found (stated up front, as the A1 feedback asked).
% =====================================================================
\section{Introduction}
\label{sec:intro}

Neural networks sometimes generalise long after they have fit their training
data, a behaviour known as \emph{grokking}~\cite{power2022grokking}. It was
first seen on small algorithmic tasks, where a network reaches perfect
training accuracy and then suddenly starts to get unseen examples right, thousands of steps later. 
This does not fit the normal picture of overfitting,
in which test performance only gets worse once the training set has been
memorised.

Liu et al.~\cite{liu2022grokking} explain grokking as representation
learning, so the network generalises once its embeddings become structured. For a
toy addition task they also make a quantitative prediction. Below a critical
training data fraction $r_c \approx 0.4$, the training data usually do not determine the right structure, and the paper
links this to no generalisation. Above, training sets increasingly determine the structure. I chose to test
this prediction because it is the sharpest claim in the paper, because each
run takes minutes on a laptop, and because the authors released their code.

This report follows the Replication Track~\cite{pineau2021reproducibility}. I
rerun the experiment behind the paper's Fig.~4 with the authors' own code,
first exactly as released and then with more seeds, and I add two ablations
on the decoder.

The theory's training-free prediction reproduces. However, the trained networks
do not follow it with the released settings. Half of the trained networks find the
structure only at $r \approx 0.76$, and in my runs having enough information in
the data was necessary but not sufficient. My ablations suggest why.
The theory predicts when the data \emph{determine} the representation, but
whether training actually \emph{finds} it depends on how fast the decoder
learns.
